% Chapitre 1 — Contexte général du projet
\chapter{Contexte général du projet}
\label{ch:contexte}

Ce chapitre situe le travail. Il présente l'organisme d'accueil, WePwn, puis le
point de départ concret du stage : un outil, \demo{}, remis en version bêta. Il pose
ensuite le problème que cet outil traite --- des agents d'intelligence artificielle
qui exécutent eux-mêmes des commandes sur la machine d'un développeur ---, fixe les
objectifs, et décrit la méthode de travail qui a permis de faire évoluer un logiciel
de sécurité sans jamais le rendre, même un instant, moins sûr.

% 1.1 — L'organisme d'accueil
\section{L'organisme d'accueil : WePwn}
\label{sec:wepwn}

\begin{tcolorbox}[enhanced,colback=wepwnclair,colframe=wepwn,boxrule=0pt,leftrule=4pt,arc=0pt,
                  left=10pt,right=10pt,top=8pt,bottom=8pt]
\begin{minipage}[c]{0.2\linewidth}
  \centering\includegraphics[width=0.85\linewidth]{figures/logos/wepwn.png}
\end{minipage}\hfill
\begin{minipage}[c]{0.76\linewidth}
  {\sffamily\Large\bfseries\color{wepwnfonce} WePwn}\par\smallskip
  {\sffamily\small Sécurité offensive · programmes de \textit{bug bounty} · divulgation
  coordonnée de vulnérabilités · outils de sécurité}
\end{minipage}
\end{tcolorbox}

\subsection{Présentation}
WePwn est une entreprise marocaine de sécurité offensive. Son activité d'origine est le
\textit{bug bounty} : une entreprise ouvre un programme, des chercheurs en sécurité
indépendants testent ses services, et chaque vulnérabilité validée est récompensée. WePwn
se présente comme la mise en relation entre des entreprises et des chercheurs, afin que
les premières soient averties des vulnérabilités identifiées chez elles~\cite{wepwn-site}.
Selon l'organisme, elle a été la première plateforme de \textit{bug bounty} au
Maroc\footnote{Affirmation de l'organisme d'accueil, reprise telle quelle.}~; les
programmes qu'elle héberge rémunèrent les chercheurs en dirhams. Un portail dédié permet
aux chercheurs de rejoindre ces programmes~\cite{wepwn-access}.

Autour de cette activité, WePwn propose un service de divulgation de vulnérabilités géré
de bout en bout~\cite{wepwn-services} : tri des signalements, évaluation de la sévérité,
validation, rapport détaillé, conseils de correction et pilotage du programme dans la
durée. Ce métier --- recevoir des attaques réelles, les reproduire, en mesurer l'impact
et accompagner la correction --- a naturellement conduit l'entreprise vers un second axe :
les \emph{pratiques} et les \emph{outils} de sécurité. \demo{} appartient à cet axe.

\subsection{Fiche signalétique}

\begin{table}[htbp]
\centering
\caption{Fiche signalétique de WePwn}
\label{tab:fiche-wepwn}
\begin{tabularx}{0.92\linewidth}{@{}L{4.2cm}Y@{}}
\toprule
Raison sociale & WePwn \\
Domaine & Cybersécurité --- sécurité offensive \\
Activités & Programmes de \textit{bug bounty} ; divulgation de vulnérabilités gérée ;
            pratiques et outils de sécurité \\
Fondateur & M. Younes Lakhouane \\
Site web & \url{https://wepwn.ma} \\
Siège & \acompleter[ville] \\
Année de création & \acompleter[année] \\
Effectif & \acompleter \\
\bottomrule
\end{tabularx}
\end{table}

\subsection{Cadre du stage}
Le stage s'est déroulé au sein de WePwn, sous l'encadrement professionnel de
M. Younes Lakhouane, fondateur de l'entreprise, et sous l'encadrement académique de
Mme K. Chetioui (ENSA de Fès). Il a commencé le \datedebut{}. Le sujet porte sur
\demo{}, un outil de WePwn publié en source ouverte sous licence
MIT~\cite{demo-cli-beta}, que l'entreprise destine aux équipes qui confient du travail
à des agents de code IA.

La double culture de l'entreprise explique la forme du travail demandé : construire
l'outil \emph{et} l'attaquer. Une partie du stage a consisté à développer et à
architecturer ; une autre, à chercher comment contourner ce qui venait d'être construit,
avec les méthodes du \textit{bug bounty} --- lire le code, fabriquer des commandes
piégées, fouiller les traces. Le chapitre~\ref{ch:securite} rend compte de cette seconde
partie.

% 1.2 — Point de départ et contribution
\section{Point de départ et contribution}
\label{sec:depart}

\subsection{La version bêta remise au début du stage}
Le stage ne part pas d'une page blanche. WePwn avait publié \demo{} en version
\betaversion{} (\betadate{})~\cite{demo-cli-beta}. Cette bêta posait déjà l'idée qui
reste au c\oe ur du projet, et qu'il faut attribuer à son point d'origine :

\begin{encadre}[L'invariant hérité de la bêta]
Une action qui modifie quelque chose doit être récupérable \emph{et} conforme au contexte
déclaré ; sinon elle est escaladée --- jamais autorisée en silence, jamais présentée à
tort comme « récupérée ».
\end{encadre}

De cet invariant découlait une posture inhabituelle pour un outil de sécurité :
\emph{récupérer plutôt que bloquer}. Quand \demo{} peut prouver qu'il a copié la cible
réelle d'une commande destructrice, il la copie et laisse l'agent continuer ; une
commande \cmd{demo\_cli undo <id>} remet les fichiers en place. Il ne bloque que ce qui
ne peut pas être récupéré (un \cmd{terraform destroy}, un \cmd{git push --force}). La bêta
s'intégrait à Claude Code et à Cursor, savait prévisualiser une requête SQL destructrice
et laissait de chaque décision un reçu chaîné.

Elle annonçait aussi, honnêtement, ses limites. Son journal des modifications le dit en
une phrase : l'analyse des commandes shell repose sur des motifs, « un arbre syntaxique
complet est sur la feuille de route ». Concrètement, le classifieur tenait en
\betaclassify{} lignes et \betaregex{} expressions régulières. Le tableau~\ref{tab:depart}
compare cette bêta à l'état du projet au commit de référence \commitref{}.

\begin{table}[htbp]
\centering
\caption{De la version bêta à l'état actuel (commit \commitref{}, \datecommitref{})}
\label{tab:depart}
\small
\begin{tabularx}{\linewidth}{@{}L{4.1cm}YY@{}}
\toprule
& \textbf{Bêta \betaversion{}} & \textbf{Aujourd'hui} \\
\midrule
Code source (Python) & $\approx$ \betalignessrc{} lignes & $\approx$ \lignessrc{} lignes \\
Tests & $\approx$ \betalignestests{} lignes, 11 fichiers & $\approx$ \lignestests{} lignes,
  \fichierstests{} fichiers \\
Lecture d'une commande & \betaregex{} expressions régulières & analyseur POSIX, simulation sur
  un système de fichiers virtuel, table de \nbregles{} règles (\nboutils{} outils) \\
Décisions & \betadecisions{} & 6 (+ SUPERVISED), canal « demander / refuser » \\
Commandes illisibles & motifs, au mieux & « non compris » : filet, point de contrôle ou
  escalade --- jamais un ALLOW \\
Pendant l'exécution & rien & garde d'appels système (Linux), montage WinFsp (Windows) \\
Surfaces distantes & quelques motifs & $\approx$ \nbsurfaces{} surfaces, chacune avec sa
  règle « ask » ou « deny » \\
Agents & Claude Code, Cursor & Claude Code (Bash, PowerShell, fichiers), Codex, Cursor,
  garde du shell \\
Fichiers de \demo{} & non protégés & plan de contrôle refusé à l'agent \\
Validation & tests unitaires & tests + exemples de règles + portail de \nbcorpus{}
  commandes, Linux et Windows \\
\bottomrule
\end{tabularx}
\end{table}

\subsection{Ce que le stage a apporté}
La mission confiée était de reprendre cette bêta, de la poursuivre et d'en reprendre
l'architecture, en menant de front le développement et le travail offensif. Le dépôt de
travail conserve le détail des \nbcommits{} commits de la dernière phase ; chacun
correspond à une partie nommée, et son message en fait le bilan.

La contribution se résume en quatre lignes :
\begin{enumerate}
  \item \textbf{Un moteur au lieu de motifs} : une analyse syntaxique unique des lignes POSIX,
    une simulation de leurs effets et une table de règles écrite comme des données
    (chapitres~\ref{ch:conception} et~\ref{ch:realisation}).
  \item \textbf{Une seule politique, un seul décideur} : ce qui était dispersé entre
    plusieurs modules est réuni dans \fichier{policy.py} et \fichier{decide.py}.
  \item \textbf{Les deux systèmes} : chaque partie est conçue et vérifiée sous Linux et sous
    Windows ; le filet d'exécution existe sur les deux.
  \item \textbf{Le travail offensif} : un audit pinçant 18 défauts, puis une suite de
    vulnérabilités (V1--V20) trouvées et corrigées, chacune avec un test qui échoue sur le
    code d'avant (chapitre~\ref{ch:securite}).
\end{enumerate}

Une part essentielle du rôle a été celle d'\emph{arbitre}. Chaque cas de politique
délicat --- faut-il autoriser un \cmd{git reset} que le \textit{reflog} sait annuler ?
demander ou refuser la suppression d'un volume Docker ? --- a fait l'objet d'une décision
écrite, numérotée et datée avant d'être codée. Les messages de commit et les tests en
gardent la trace : « ruling (c) », « rulings M1--M5 », « Nizar's table ».

% 1.3 — Problématique
\section{Problématique : des agents qui exécutent des commandes}
\label{sec:problematique}

\subsection{Un agent n'écrit plus seulement du code, il l'exécute}
Un agent de code IA comme Claude Code ou Codex ne se contente pas de proposer du texte.
Il lit le projet, modifie des fichiers, et surtout lance des commandes dans un shell ---
\cmd{bash} sous Linux, PowerShell sous Windows --- pour compiler, tester, nettoyer, migrer
une base, déployer. Chaque commande s'exécute avec les droits du développeur, sur sa
machine, souvent à côté de données qui n'existent nulle part ailleurs. L'agent se trompe
parfois ; une erreur dans une commande destructrice n'est pas une suggestion maladroite,
c'est une perte de données.

\subsection{Trois incidents publics}
Le risque n'est pas théorique. Trois cas documentés illustrent trois formes du problème.

\begin{description}[style=nextline,leftmargin=0.6cm,font=\sffamily\bfseries]
  \item[Une base de production supprimée (Replit, juillet 2025)]
    Au cours d'une expérience de plusieurs jours, l'agent de Replit a exécuté des commandes
    qui ont effacé la base de données de production d'un utilisateur, alors qu'il avait reçu
    la consigne de ne rien modifier sans accord explicite ; il a ensuite affirmé à tort que
    la restauration était impossible~\cite{oecd-replit,eweek-replit}.
    \emph{Leçon : une consigne n'est pas une protection, et l'agent n'est pas une source
    fiable sur ce qui peut être récupéré.}
  \item[Des fichiers écrasés par des déplacements (Gemini CLI, juillet 2025)]
    L'agent a cru avoir créé un dossier dont la création avait échoué, puis y a « déplacé »
    les fichiers de l'utilisateur ; chaque déplacement vers un nom inexistant a écrasé le
    précédent, et tous les fichiers sauf un ont été perdus~\cite{aiid-gemini,gh-gemini}.
    \emph{Leçon : aucune commande de la séquence ne « ressemble » à une suppression ; seul
    un raisonnement sur l'état du système de fichiers révèle la perte.}
  \item[Des fichiers supprimés en voulant les compter (Claude Code, juillet 2026)]
    Chargé de compter des fichiers de rapport, un agent a lancé une commande qui les a
    supprimés ; l'utilisateur en a récupéré une partie avec un outil de récupération de
    données~\cite{gh-claude-76626}. La bêta de \demo{} cite ce cas comme motivation.
    \emph{Leçon : demander une confirmation ne suffit pas quand la personne fait confiance
    à l'intention annoncée ; ce qui manquait, c'était une copie prise avant.}
\end{description}

\subsection{Pourquoi le problème est difficile}
Juger une commande \emph{avant} qu'elle s'exécute paraît simple : repérer \cmd{rm}, et se
méfier. En réalité le langage du shell rend cette lecture trompeuse. Le
tableau~\ref{tab:regex-pieges} montre quatre lignes, tirées des tests du projet, où une
lecture naïve --- par motifs, ou par un analyseur trop confiant --- se trompe dans un sens
ou dans l'autre.

\begin{table}[htbp]
\centering
\caption{Lignes où une lecture naïve se trompe (extraits de \fichier{tests/test\_command\_model.py} et des commits I1a, W0)}
\label{tab:regex-pieges}
\small
\begin{tabularx}{\linewidth}{@{}L{5.2cm}YY@{}}
\toprule
\textbf{Ligne} & \textbf{Lecture naïve} & \textbf{Ce que fait le shell} \\
\midrule
\cmd{ls \# ; rm -rf x} & un \cmd{rm -rf} & rien : tout ce qui suit \cmd{\#} est un commentaire \\
\cmd{echo \$(date; rm -rf x) >/dev/null} & un \cmd{echo} & exécute \cmd{rm -rf x} dans la substitution \\
\cmd{bash -o pipefail -c "rm app.db"} & une option inconnue, rien à juger & supprime \cmd{app.db} \\
\cmd{X=-rf; rm \$X src} & un \cmd{rm} sans option & un \cmd{rm -rf src} \\
\bottomrule
\end{tabularx}
\end{table}

À cela s'ajoutent les enveloppes (\cmd{sudo}, \cmd{xargs}, \cmd{timeout}), les commandes
lancées par du code que l'outil ne lit pas (\cmd{python script.py}), les effets à distance
(une base, un dépôt git distant, un bucket) qu'aucune copie locale ne restaure, et un
second langage, PowerShell, où les arguments se lient aux paramètres par position.

\subsection{Énoncé du problème}
\begin{encadre}[Problématique]
Comment, avant qu'une commande écrite par un agent ne s'exécute, en prédire l'effet assez
exactement pour laisser passer ce qui est sans danger, rendre récupérable ce qui peut
l'être, et arrêter ou soumettre à une personne le reste --- sans jamais revendiquer une
protection qu'on n'a pas prise, et de la même manière sous Linux et sous Windows ?
\end{encadre}

% 1.4 — Objectifs
\section{Objectifs du projet}
\label{sec:objectifs}

Les objectifs découlent de la problématique ; chacun est formulé de façon à pouvoir être
vérifié dans le code et les tests.

\begin{table}[htbp]
\centering
\caption{Objectifs du stage et critère de vérification de chacun}
\label{tab:objectifs}
\small
\begin{tabularx}{\linewidth}{@{}L{0.7cm}L{5.6cm}Y@{}}
\toprule
& \textbf{Objectif} & \textbf{Comment on le vérifie} \\
\midrule
O1 & \textbf{Prédire} l'effet d'une ligne à partir de sa structure, pas de son texte. &
  Chaque exemple de la table de règles (\nbexemples{}) est un test qui compare l'impact
  prédit à l'impact attendu. \\
O2 & \textbf{Ne jamais autoriser ce qui n'est pas lu.} Une ligne illisible n'obtient pas
  ALLOW. & Tests M2.3a : une ligne « non comprise » va au filet, à un point de contrôle,
  sinon elle est escaladée. \\
O3 & \textbf{Ne jamais revendiquer une récupération non prise.} &
  Tests des captures partielles (tout ou rien), de V17 (point de contrôle hors projet). \\
O4 & \textbf{Voir pendant l'exécution} ce que la lecture ne peut pas prédire. &
  Garde d'appels système (Linux) et montage WinFsp (Windows) testés sur les deux systèmes. \\
O5 & \textbf{Traiter le distant honnêtement} : distinguer « demander à la personne » et
  « refuser ». & Table \fichier{policy.ESCALATION\_CHANNEL} ; une surface sans ligne est
  refusée. \\
O6 & \textbf{Les deux systèmes} : Linux et Windows, avec les mêmes garanties. &
  Suite complète passée sur les deux systèmes avant chaque atterrissage. \\
O7 & \textbf{Ne rien rendre plus permissif sans le vouloir.} &
  Portail de décision : \nbcorpus{} commandes comparées avant et après chaque partie. \\
O8 & \textbf{Résister à l'agent} : il ne doit pas pouvoir désarmer l'outil. &
  Plan de contrôle refusé à toutes les portes ; tâche Windows durcie (W-SEC). \\
\bottomrule
\end{tabularx}
\end{table}

Un objectif est volontairement \emph{absent} : faire de \demo{} une barrière contre un
adversaire déterminé. L'outil protège contre les \emph{erreurs} d'un agent coopératif ;
ce modèle de menace, hérité de la bêta et affiché par \cmd{demo\_cli doctor}, est discuté
à la section~\ref{sec:menace}.

% 1.5 — Méthodologie
\section{Méthodologie de travail}
\label{sec:methodologie}

\subsection{Le risque propre à un outil de sécurité}
Faire évoluer un outil de sécurité comporte un risque particulier : une modification qui
corrige un cas peut, sans bruit, en autoriser un autre. Un test « vert » ne le révèle pas
si aucun test ne décrivait le cas devenu permissif. Un audit mené au cours du stage l'a montré :
la suite passait avec 1856 tests réussis alors que 18 défauts confirmés restaient
atteignables, parce qu'aucun test ne décrivait ces comportements. La méthode adoptée
répond directement à ce constat.

\subsection{Une partie à la fois}
Le travail est découpé en \emph{parties} courtes et nommées (A, B, C1, SYS1, M2.2b3,
W-SEC\ldots). Chaque partie suit le même cycle (figure~\ref{fig:methode}) :

\begin{enumerate}
  \item \textbf{Plan} : périmètre, cas à couvrir, risques, ce qui ne doit pas bouger.
  \item \textbf{Arbitrages} : chaque question de politique est tranchée \emph{par écrit}
    avant d'être codée, avec un identifiant et une date (« ruling R2 : un envoi qui peut
    remplacer un objet distant est escaladé »). Le code et ses tests citent l'arbitrage.
  \item \textbf{Prototype} : le code, et des tests qui \emph{échouent sur le code d'avant}.
    Un test qui passerait aussi sur l'ancien code ne prouve rien sur la correction.
  \item \textbf{Revue} : relecture adverse et sondes --- des commandes fabriquées pour
    piéger la nouvelle logique, lancées dans un projet jetable.
  \item \textbf{Atterrissage} : un commit dont le message est un bilan : ce qui change, les
    arbitrages appliqués, les tests ajoutés, et le résultat du portail de décision.
\end{enumerate}

\begin{figure}[htbp]
\centering
% Figure 6 — La méthode : une partie à la fois, et les portes qu'elle doit franchir.
\begin{tikzpicture}[
    etape/.style={accent,text width=24mm,minimum height=15mm,font=\sffamily\small},
    porte/.style={boite,draw=encre,text width=30mm,font=\sffamily\scriptsize,minimum height=9mm},
  ]
  \node[etape] (p) at (0,0) {Plan\sous{périmètre, cas, risques}};
  \node[etape,right=6mm of p] (r) {Arbitrages\sous{décisions écrites et datées (R1, M3\ldots)}};
  \node[etape,right=6mm of r] (pr) {Prototype\sous{code + tests qui échouent sur l'ancien code}};
  \node[etape,right=6mm of pr] (rv) {Revue\sous{lecture adverse, sondes}};
  \node[etape,right=6mm of rv,fill=wepwn!35] (l) {Atterrissage\sous{un commit, message = bilan}};
  \foreach \a/\b in {p/r,r/pr,pr/rv,rv/l} \draw[fleche] (\a) -- (\b);
  \draw[fleche,densely dashed] (rv.north) -- ++(0,0.35) -| node[etiquette,pos=0.25,above]{correction} (pr.north);

  % Portes
  \node[porte] (g1) at (2.0,-2.8) {Tests de la partie\sous{rouges avant, verts après}};
  \node[porte,right=4mm of g1] (g2) {Suite complète\sous{Linux \emph{et} Windows}};
  \node[porte,right=4mm of g2] (g3) {Portail de décision\sous{\nbcorpus{} commandes, avant/après}};
  \node[porte,right=4mm of g3] (g4) {Aucune décision plus lâche\sous{sans arbitrage explicite}};
  \node[zone,fit=(g1)(g4),label={[etiquette]below:portes à franchir avant l'atterrissage}] (gz) {};
  \draw[fleche] (gz.north -| l.south) -- (l.south);

  % Boucle de parties
  \node[etiquette,text width=120mm,anchor=north] at ([yshift=-7mm]gz.south) {Parties successives (extrait) :
     A $\cdot$ B $\cdot$ C1 $\cdot$ C2a $\cdot$ SYS1 $\cdot$ S2+M1 $\cdot$ M2.0--M2.1 $\cdot$ T $\cdot$ T5 $\cdot$ B0 $\cdot$ B1 $\cdot$ I1 $\cdot$ W0 $\cdot$
     M2.2a--f $\cdot$ ASK $\cdot$ W-SEC $\cdot$ V17 $\cdot$ M2.3a $\cdot$ M2.3b $\cdot$ V22 $\cdot$ V18 $\cdot$ C3};
\end{tikzpicture}

\caption{La méthode : le cycle d'une partie et les portes à franchir avant l'atterrissage}
\label{fig:methode}
\end{figure}

\subsection{Les portes d'un atterrissage}
\begin{description}[style=nextline,leftmargin=0.6cm,font=\sffamily\bfseries]
  \item[Les tests de la partie] Rouges sur le code d'avant, verts après.
  \item[La suite complète, sous Linux et sous Windows] Les deux systèmes n'ont pas les mêmes
    chemins (\cmd{C:\textbackslash}, \cmd{/c/} sous Git Bash), ni les mêmes appels système, ni le
    même shell : un test vert sous Linux ne dit rien de Windows.
  \item[Le portail de décision] Un corpus de \nbcorpus{} commandes réelles et fabriquées est
    évalué avant et après la modification ; toute ligne dont la décision, la règle, la
    cible ou la capture change est listée et doit être justifiée. Les messages de commit en
    rendent compte (« the 922 gate changes 124 reviewed decisions: 8 fixed misses,
    18 newly captured, 86 newly allowed »).
  \item[Aucune décision plus permissive sans arbitrage] Une ligne qui passe d'ESCALATE à ALLOW
    sans arbitrage explicite bloque l'atterrissage, même si tous les tests passent.
\end{description}

\begin{pourquoi}[Pourquoi un portail de décision en plus des tests ?]
Les tests vérifient ce qu'on a pensé à écrire. Le portail vérifie \emph{tout le reste} :
il ne sait pas quelle décision est juste, mais il voit chaque décision qui change. C'est
lui qui transforme « la suite passe » en « rien n'a bougé sans qu'on le sache ».
\end{pourquoi}

\subsection{Le message de commit comme compte rendu}
Chaque atterrissage laisse un compte rendu lisible directement dans l'historique.
L'extrait suivant est réel (les lignes « [...] » sont coupées) :

\begin{terminal}[git show -s 3c94414]
$ git show -s --format=%B 3c94414
M2.3a: lines the engine does not fully read leave the old text rules

- runners (python, node, npm, make, ./x, source, awk programs that run or
  write ...) go to the runtime net, else a checkpoint, else ALLOW with a
  note that nothing was captured (ruling M1)
- a line the engine cannot read goes to the net, else a checkpoint, else
  ESCALATE (M3) - unread SQL such as psql -c "$(cat q.sql)" was ALLOWED
  (V19); where the old rules flag the line they keep their capture or
  escalation, never an ALLOW
[...]
- engine rules for gzip / bzip2 / xz, sort -o, iconv -o, patch,
  awk -i inplace, 7z: they deleted or overwrote files unseen (M4, V20)
[...]
- tests/test_m23a.py (17); row fixtures of earlier parts: every moved row
  stricter; gate: 6 rows move, reference re-recorded
\end{terminal}

% 1.6 — Planning
\section{Planning}
\label{sec:planning}

\subsection{Déroulement du stage}
Le stage a débuté le \datedebut{} et s'est terminé le \datefin.
\acompleter[phases du stage, de la prise en main de la bêta à la fin : à confirmer pour le
diagramme de Gantt d'ensemble]

\subsection{Chronologie des parties documentées}
L'historique du dépôt de travail permet de dater précisément les parties livrées entre le
21 septembre et le 9 octobre 2026. La figure~\ref{fig:gantt} les regroupe en cinq phases :
la reprise et l'audit ; un moteur qui prédit sans encore décider (ses prédictions sont
seulement écrites dans les reçus, pour être comparées aux décisions existantes) ; le moteur
qui décide ; l'extension famille par famille de la table de règles ; enfin la sécurité et
la sortie de l'ancien chemin par expressions régulières. La partie C3 (le moteur
PowerShell) est \statutCthree{}.

\begin{figure}[htbp]
\centering
% Figure 7 — Chronologie des parties livrées, d'après l'historique git (21/09 -> 09/10/2026).
% Une barre = les jours où les commits de la partie ont atterri.
\begin{ganttchart}[
    x unit=5.2mm, y unit chart=5.2mm, y unit title=5.5mm,
    time slot format=isodate, vgrid={*1{draw=trait,dotted}},
    title/.style={fill=grisclair,draw=trait},
    title label font=\sffamily\scriptsize,
    bar/.style={fill=wepwn!70,draw=wepwnfonce},
    bar height=0.55, bar label font=\sffamily\scriptsize,
    group/.style={fill=encre!80,draw=encre}, group label font=\sffamily\scriptsize\bfseries,
    group right shift=0, group top shift=0.55, group height=0.25, group peaks height=0.12,
    milestone/.style={fill=dASK,draw=dASK!70!black}, milestone label font=\sffamily\scriptsize,
    link/.style={-latex,draw=gris},
  ]{2026-09-21}{2026-10-11}
  \gantttitle{Jours du cycle de développement documenté}{21} \\
  \gantttitlelist{1,...,21}{1} \\
  \ganttgroup{Reprise et audit}{2026-09-21}{2026-09-27} \\
  \ganttbar{Import v1.7.0, outillage}{2026-09-21}{2026-09-22} \\
  \ganttbar{Audit : 18 défauts}{2026-09-26}{2026-09-27} \\
  \ganttbar{Modèle de commande}{2026-09-27}{2026-09-27} \\
  \ganttgroup{Moteur (prédiction seule)}{2026-09-28}{2026-09-29} \\
  \ganttbar{A, B, C1, C2a}{2026-09-28}{2026-09-28} \\
  \ganttbar{SYS1 : filets reliés}{2026-09-29}{2026-09-29} \\
  \ganttgroup{Le moteur décide}{2026-10-02}{2026-10-06} \\
  \ganttbar{S2+M1, M2.0+M2.1}{2026-10-02}{2026-10-02} \\
  \ganttbar{T, T5 : une politique}{2026-10-03}{2026-10-04} \\
  \ganttbar{B0, B1, B1b}{2026-10-05}{2026-10-06} \\
  \ganttbar{I1a--c, W0}{2026-10-06}{2026-10-06} \\
  \ganttgroup{Familles d'outils (M2.2)}{2026-10-06}{2026-10-08} \\
  \ganttbar{git, infra, cloud}{2026-10-06}{2026-10-07} \\
  \ganttbar{conteneurs, données, ASK}{2026-10-07}{2026-10-07} \\
  \ganttbar{envois, PaaS, SaaS}{2026-10-07}{2026-10-07} \\
  \ganttbar{système, LATER, préfiltre}{2026-10-08}{2026-10-08} \\
  \ganttgroup{Sécurité et sortie de l'ancien chemin}{2026-10-08}{2026-10-09} \\
  \ganttbar{W-SEC, V17}{2026-10-08}{2026-10-08} \\
  \ganttbar{M2.3a, M2.3b, V22, V18}{2026-10-09}{2026-10-09} \\
  \ganttbar[bar/.style={fill=white,draw=wepwnfonce,dashed}]{C3 (PowerShell)}{2026-10-10}{2026-10-11}
\end{ganttchart}

\caption{Chronologie des parties livrées, d'après l'historique git (21 septembre -- 9 octobre 2026)}
\label{fig:gantt}
\end{figure}

\begin{pourquoi}[Pourquoi un moteur « en ombre » avant de le laisser décider ?]
Les parties B à C2a ont livré un moteur qui \emph{ne décidait rien} : il écrivait sa
prédiction dans chaque reçu, et le portail vérifiait que zéro décision changeait. Pendant
ce temps, chaque divergence entre la prédiction et la décision existante était un
signal : soit un défaut du moteur, soit un défaut de l'ancien chemin. Le moteur n'a pris
la main (partie S2+M1) qu'une fois ces divergences comprises, et ce changement-là a été
revu ligne par ligne sur le portail.
\end{pourquoi}


\section*{Conclusion du chapitre}
Le stage part d'un outil qui existait, qui fonctionnait sur les cas simples et qui
annonçait lui-même sa limite principale : il lisait les commandes avec des expressions
régulières. Le problème est réel et documenté par des incidents publics ; les objectifs
sont mesurables ; la méthode impose qu'aucune décision ne devienne plus permissive sans
qu'on l'ait voulu. Le chapitre suivant examine ce qui existe déjà pour protéger un
développeur de son agent, et pourquoi ces protections ne suffisent pas.
